\section{Group Action on a Set}

For a group action, $gx$ denotes the result of applying $g$ to $x$; it need
not be multiplication in a group containing $x$. The axioms are $ex=x$ and
$(gh)x=g(hx)$. The orbit $Gx$ consists of all possible images of $x$, the
stabilizer $G_x$ consists of the group elements fixing $x$, and $X_g$ is
the set of points fixed by one specified $g$. Keep these three sets distinct.
In the classification exercises, a transitive set is understood to be
nonempty, as in the book's examples. If the empty set is admitted, the
quantified transitivity condition holds vacuously there and adds one empty case.

\exercise{16.1}{All fixed sets in the square action}
\idea In each row of Table 16.10, retain exactly the column labels that the
row leaves unchanged. Write $X$ for the full seventeen-element set.
\begin{align*}
X_{\rho_0}&=X,&X_{\rho_1}&=\{C\},&X_{\rho_3}&=\{C\},\\
X_{\rho_2}&=\{m_1,m_2,d_1,d_2,C\},\\
X_{\mu_1}&=\{s_1,s_3,m_1,m_2,C,P_1,P_3\},\\
X_{\mu_2}&=\{s_2,s_4,m_1,m_2,C,P_2,P_4\},\\
X_{\delta_1}&=\{2,4,d_1,d_2,C\},\\
X_{\delta_2}&=\{1,3,d_1,d_2,C\}.
\end{align*}
For example, a half-turn exchanges opposite vertices, sides, and midpoints,
but leaves each axis and each diagonal as a set. An object in $X$ can be a
whole line or side, so ``fixed'' does not require each point of that object
to be fixed individually.

\exercise{16.2}{Every stabilizer in the square action}
\idea Read down a column and collect the row labels that reproduce its heading.
Let $e=\rho_0$. The answers can be grouped by identical stabilizers:
\[
\begin{array}{c|l}
 x&G_x\\\hline
1,3&\{e,\delta_2\}\\
2,4&\{e,\delta_1\}\\
s_1,s_3,P_1,P_3&\{e,\mu_1\}\\
s_2,s_4,P_2,P_4&\{e,\mu_2\}\\
m_1,m_2&\{e,\rho_2,\mu_1,\mu_2\}\\
d_1,d_2&\{e,\rho_2,\delta_1,\delta_2\}\\
C&D_4
\end{array}
\]
For example, vertex $1$ is fixed by the identity and by reflection in its
own diagonal $d_2$, and by no other square symmetry. The orbit--stabilizer
check gives orbit size $8/2=4$, as expected for the vertices.

\exercise{16.3}{The six orbits of the square action}
Starting from a vertex reaches every vertex and no other kind of object.
The same observation for each kind of object gives
\[
\{1,2,3,4\},\quad\{s_1,s_2,s_3,s_4\},\quad
\{m_1,m_2\},\quad\{d_1,d_2\},\quad\{C\},\quad
\{P_1,P_2,P_3,P_4\}.
\]
They are disjoint, their sizes sum to $17$, and each is closed under all
square symmetries. This verifies that the list exhausts $X$.

\exercise{16.4}{Correcting faithful action}
The quantifier ``for every $x$'' is essential. An action is faithful if
\[
\bigl(gx=x\text{ for every }x\in X\bigr)\ \Longrightarrow\ g=e.
\]
Equivalently, the homomorphism $G\to S_X$ determined by the action is injective.
Requiring this implication for a single fixed point would mean that all
point stabilizers are trivial, which is the stronger condition called a
free action. A reflection of a square fixes some vertices without acting
as the identity on all vertices.

\exercise{16.5}{Correcting transitive action}
For every $x,y\in X$ there must exist $g\in G$ such that $gx=y$. The group
element may depend on both points. One fixed $g$ cannot send a single $x$
to several different outputs. For a nonempty $X$, an equivalent formulation
is: the orbit of any one chosen point is all of $X$.

\exercise{16.6}{Subsets preserved by an action}
\answer{The sub-$G$-sets are precisely the unions of whole $G$-orbits.}
If $S$ is a sub-$G$-set and $s\in S$, closure gives $Gs\subseteq S$.
Since $s=es$, every $s$ is in its own orbit, so $S=\bigcup_{s\in S}Gs$.
Conversely, if $S$ is a union of orbits and $s\in Gx\subseteq S$, then
$s=hx$ for some $h$. For any $g$, $gs=(gh)x\in Gx\subseteq S$.
The empty union also yields the empty sub-$G$-set.

\exercise{16.7}{Transitivity in terms of orbits}
For nonempty $X$, transitivity is equivalent to having exactly one orbit.
If every $y$ equals $gx$ for some $g$, then $Gx=X$. Conversely, when
$Gx=X$, any $y=hx$ and $z=kx$ satisfy $(kh^{-1})y=z$.
If empty sets are allowed under the literal quantified definition, the
corresponding characterization is ``at most one orbit.''

\exercise{16.8}{True or false: action properties}
\begin{enumerate}[label=(\alph*)]
\item False. A group can act on any set; a three-element set with the trivial
$\Z_2$ action need not come equipped with a group operation.
\item True. This is the axiom $ex=x$.
\item False. In the trivial action every $g$ fixes every point; faithfulness
is an additional condition.
\item True. Apply $g^{-1}$ to $gx_1=gx_2$ and use the action axioms to get
$x_1=x_2$.
\item False. In the trivial action of a nontrivial group all group elements
send a chosen $x$ to the same point.
\item True. It is preserved by $G$, and if $hx,kx$ are two points in the
orbit then $kh^{-1}$ sends the first to the second.
\item True. Restrict the action map to $H\times X$; the same two axioms hold.
\item False. Take $H=\{e\}$ inside a group that acts transitively on more than
one point. Its orbits are singletons, while the original orbit is larger.
\item True. The inverse of the map $x\mapsto gx$ is $x\mapsto g^{-1}x$.
\item True. Orbit--stabilizer gives $|Gx|=[G:G_x]=|G|/|G_x|$ when $G$ is finite.
\end{enumerate}

\exercise{16.9}{Isomorphic orbits and other subactions}
\begin{enumerate}[label=(\alph*)]
\item The side orbit and midpoint orbit are isomorphic by $s_i\mapsto P_i$.
Every square symmetry sends the midpoint of a side to the midpoint of its
image side, so this bijection commutes with the action.
\item The reflection $\mu_1$ fixes two sides, $s_1,s_3$, but no vertices.
An equivariant bijection must carry the fixed points of any $g$ bijectively
onto the fixed points of the same $g$. Thus no such bijection can identify
the vertex and side orbits.
\item No. The question permits sub-$D_4$-sets that are unions of orbits.
For instance $\{s_1,s_2,s_3,s_4,C\}$ and
$\{P_1,P_2,P_3,P_4,C\}$ are different isomorphic sub-$D_4$-sets: extend the
bijection in part (a) by $C\mapsto C$. Among individual orbits, the only
isomorphic pair is indeed the side and midpoint pair: the other four-point
orbit is excluded by part (b), while $\mu_1$ fixes both axes and exchanges
the diagonals, distinguishing the two two-point orbits.
\end{enumerate}

\exercise{16.10}{Faithfulness of the square action}
\begin{enumerate}[label=(\alph*)]
\item Yes. A symmetry fixing all seventeen objects fixes all four vertices.
Only the identity symmetry does that.
\item The vertex, side, and midpoint orbits each give faithful actions:
inspection of their columns in Table 16.10 shows that no nonidentity row
fixes the entire orbit. The axis orbit and diagonal orbit are not faithful,
since the nonidentity half-turn $\rho_2$ fixes both members in either case.
The center orbit is not faithful, since every symmetry fixes $C$.
\end{enumerate}

\exercise{16.11}{Faithfulness and distinct permutations}
\goal Prove that faithfulness says precisely that different group elements
induce different maps on $X$.
Suppose first that the action is faithful and $gx=hx$ for every $x$.
Applying $h^{-1}$ gives $(h^{-1}g)x=x$ for all $x$, so $h^{-1}g=e$ and $g=h$.
Conversely, if distinct group elements have distinct actions and $gx=x$ for
every $x$, then $g$ and $e$ induce the same map. Hence $g=e$, proving
faithfulness. Equality of actions here means equality on \emph{every} point,
not merely agreement at one point.

\exercise{16.12}{Pointwise stabilizers of subsets}
Let $G_Y=\{g:gy=y\text{ for every }y\in Y\}$. The identity belongs to $G_Y$.
If $g,h\in G_Y$, then for every $y\in Y$,
$(gh)y=g(hy)=gy=y$, so $gh\in G_Y$. If $g\in G_Y$, applying $g^{-1}$ to
$gy=y$ gives $g^{-1}y=y$, so $g^{-1}\in G_Y$.
Thus $G_Y$ is a subgroup. Equivalently $G_Y=\bigcap_{y\in Y}G_y$.
When $Y=\varnothing$, all conditions are vacuous and $G_Y=G$.

\exercise{16.13}{The real rotation action}
\begin{enumerate}[label=(\alph*)]
\item For $\theta\in(\R,+)$ use
$R_\theta=\begin{bmatrix}\cos\theta&-\sin\theta\\\sin\theta&\cos\theta\end{bmatrix}$.
We have $R_0=I$ and the angle-addition identities give
$R_{\theta+\phi}=R_\theta R_\phi$. These are the two action axioms.
\item Rotation preserves distance to the origin. Conversely every point
at distance $|P|>0$ has an angle differing from that of $P$ by some real
$\theta$. Hence the orbit is the entire circle centered at the origin
through $P$.
\item A nonzero point returns to itself exactly after an integer multiple
of a full turn, so $G_P=2\pi\Z$ as a subgroup of $(\R,+)$.
In particular, this real-number action is not faithful: $2\pi$ acts trivially.
\end{enumerate}

\exercise{16.14}{Unions of disjoint actions}
\begin{enumerate}[label=(\alph*)]
\item If the $X_i$ are disjoint, every $x$ in their union belongs to exactly
one $X_i$. Define $gx$ using that component's given action. This is
well defined, stays in $X_i$, and satisfies $ex=x$ and $(gh)x=g(hx)$ because
the component action does.
\item Every $x\in X$ belongs to $Gx$. Two orbits meeting at a point are
equal: if $gx=hy$, then $x=g^{-1}hy$, so every image of $x$ is an image of
$y$, and conversely. Thus the distinct orbits partition $X$, and their
union, with the component actions, is exactly the original action.
\end{enumerate}

\exercise{16.15}{A transitive action is a coset action}
Put $H=G_{x_0}$ and define $\phi(gx_0)=gH$; transitivity provides such a $g$
for each point.
\step{Prove the definition does not depend on the choice of $g$.}
If $gx_0=hx_0$, then $(h^{-1}g)x_0=x_0$, so $h^{-1}g\in H$ and $gH=hH$.
The same equivalences in reverse show that $gH=hH$ implies $gx_0=hx_0$.
Thus $\phi$ is both well defined and injective.
\step{Check surjectivity and the action.}
Every coset $gH$ is the image of $gx_0$, so $\phi$ is onto. For $k\in G$,
\[
\phi(k(gx_0))=\phi((kg)x_0)=(kg)H=k\phi(gx_0).
\]
This proves the required $G$-set isomorphism.

\exercise{16.16}{Decomposing an arbitrary action}
First partition $X$ into its distinct orbits $O_i$. Choose a point $x_i$ in
each nonempty orbit and let $H_i=G_{x_i}$. Exercise 16.15 supplies an
isomorphism $O_i\to G/H_i$, where $G/H_i$ here denotes the \emph{set of left
cosets}; $H_i$ need not be normal. Tag the coset set by $i$ and let
\[
g\cdot(kH_i,i)=(gkH_i,i).
\]
The map taking $x\in O_i$ to $(\phi_i(x),i)$ is bijective, because the orbits
partition $X$ and each $\phi_i$ is bijective. It is equivariant in each
component and hence globally. The tags prevent cosets from different
components from being confused. For an arbitrary infinite family, choosing
all the representatives uses the usual choice convention; finite actions
need no such qualification.

\exercise{16.17}{When two coset actions are isomorphic}
\begin{enumerate}[label=(\alph*)]
\item Let $H=G_{x_0}$. An element $k$ fixes $g_0x_0$ exactly when
$(g_0^{-1}kg_0)x_0=x_0$, equivalently $g_0^{-1}kg_0\in H$.
Consequently $G_{g_0x_0}=g_0Hg_0^{-1}$.
\item The conjecture is that $G/H$ and $G/K$ are isomorphic as $G$-sets
exactly when $H$ and $K$ are conjugate subgroups.
\item If $K=tHt^{-1}$, define $\phi(gH)=gt^{-1}K$.
Replacing $g$ by $gh$ for $h\in H$ does not change the answer, because
$ght^{-1}K=gt^{-1}(tht^{-1})K=gt^{-1}K$.
Its inverse sends $uK$ to $utH$, and left multiplication commutes with
$\phi$, proving sufficiency. Conversely, an equivariant bijection sends
$H$ to some $aK$ and preserves stabilizers. The stabilizer of $H$ is $H$,
while that of $aK$ is $aKa^{-1}$. Thus $H=aKa^{-1}$, proving necessity.
\end{enumerate}

\exercise{16.18}{All transitive actions of the group of order four}
\idea A transitive action is a coset action, and isomorphism types correspond
to conjugacy classes of subgroups. Since $\Z_4$ is abelian, conjugation does
not identify different subgroups. Its subgroups are $\Z_4$, $\{0,2\}$, and
$\{0\}$, giving exactly three types, with $1,2,4$ points respectively.
Their action tables are
\[
\begin{array}{c|c} &a\\\hline0&a\\1&a\\2&a\\3&a\end{array}
\qquad
\begin{array}{c|cc}&a&b\\\hline0&a&b\\1&b&a\\2&a&b\\3&b&a\end{array}
\qquad
\begin{array}{c|cccc}&a&b&c&d\\\hline
0&a&b&c&d\\1&b&c&d&a\\2&c&d&a&b\\3&d&a&b&c
\end{array}.
\]
In the four-point table, $a,b,c,d$ represent $0,1,2,3$; in the two-point
table they represent the two cosets of $\{0,2\}$.

\exercise{16.19}{All transitive actions of the group of order six}
There is one subgroup of each order $1,2,3,6$ in $\Z_6$ and no others,
since every subgroup of a cyclic group is cyclic. Conjugation is trivial.
The four isomorphism types have sizes $6,3,2,1$.
For the first three smaller sizes the tables are
\[
\begin{array}{c|c}&a\\\hline0&a\\1&a\\2&a\\3&a\\4&a\\5&a\end{array}
\qquad
\begin{array}{c|cc}&a&b\\\hline0&a&b\\1&b&a\\2&a&b\\3&b&a\\4&a&b\\5&b&a\end{array}
\qquad
\begin{array}{c|ccc}&a&b&c\\\hline
0&a&b&c\\1&b&c&a\\2&c&a&b\\3&a&b&c\\4&b&c&a\\5&c&a&b
\end{array}.
\]
The six-point regular action is
\[
\begin{array}{c|cccccc}&a&b&c&d&e&f\\\hline
0&a&b&c&d&e&f\\1&b&c&d&e&f&a\\2&c&d&e&f&a&b\\
3&d&e&f&a&b&c\\4&e&f&a&b&c&d\\5&f&a&b&c&d&e
\end{array}.
\]
Each row adds its group element to the residue represented by the column,
modulo the number of points. Since that number divides $6$, this defines
an action of $\Z_6$ in every case.

\exercise{16.20}{All transitive actions of the symmetric group on three letters}
Write $r=(1,2,3)$, $r^2=(1,3,2)$, $s=(2,3)$, $t=(1,3)$, and $u=(1,2)$.
The subgroups are $S_3$, $A_3$, the three conjugate order-two subgroups,
and the identity subgroup. Hence there are four transitive isomorphism
types, with $1,2,3,6$ points. The one-point action is trivial; the two-point
action records even or odd parity; the three-point action is the natural
permutation action; and the six-point action is left multiplication on $S_3$.
Here are all their action tables, with the requested row order:
\[
\begin{array}{c|c}&a\\\hline e&a\\r&a\\r^2&a\\s&a\\t&a\\u&a\end{array}
\quad
\begin{array}{c|cc}&a&b\\\hline e&a&b\\r&a&b\\r^2&a&b\\s&b&a\\t&b&a\\u&b&a\end{array}
\quad
\begin{array}{c|ccc}&a&b&c\\\hline
e&a&b&c\\r&b&c&a\\r^2&c&a&b\\s&a&c&b\\t&c&b&a\\u&b&a&c
\end{array},
\]
\[
\begin{array}{c|cccccc}&a&b&c&d&e&f\\\hline
e&a&b&c&d&e&f\\r&b&c&a&f&d&e\\r^2&c&a&b&e&f&d\\
s&d&e&f&a&b&c\\t&e&f&d&c&a&b\\u&f&d&e&b&c&a
\end{array}.
\]
In the last table the point labels $a,b,c,d,e,f$ represent the group
elements $e,r,r^2,s,t,u$ respectively. For example, $rs=u$, so row $r$,
column $d$ has entry $f$. The point label $e$ in that table is distinct
from the row label for the group identity.
